\documentclass[12pt]{article}
\usepackage{amsmath,amssymb,amsfonts}
\usepackage{tikz-cd}
\usepackage{geometry}
\geometry{margin=1in}

\title{Morphisms in the Opposite Fibration via Spans}
\author{Compiled for Siya}
\date{}

\begin{document}
	
	\maketitle
	
	\section*{Setup}
	
	Let \( P : \mathbf{E} \to \mathbf{B} \) be a fibration. The opposite fibration is defined as:
	\[
	P^{\mathrm{op}} : \mathbf{E}^{\mathrm{op}} \to \mathbf{B}^{\mathrm{op}}
	\]
	We aim to describe morphisms in \( P^{\mathrm{op}} \) using spans in the original category \( \mathbf{E} \).
	
	\section*{Morphisms in \( P^{\mathrm{op}} \)}
	
	Given:
	\begin{itemize}
		\item An object \( Y \in \mathbf{E}_J \)
		\item An object \( X \in \mathbf{E}_I \)
		\item A morphism \( u : J \to I \) in \( \mathbf{B}^{\mathrm{op}} \), i.e., \( u : I \to J \) in \( \mathbf{B} \)
	\end{itemize}
	
	A morphism in \( P^{\mathrm{op}} \) from \( Y \) to \( X \) over \( u \) is encoded by a span in \( \mathbf{E} \):
	\[
	X \xleftarrow{\text{vertical}} e \xrightarrow{\text{cartesian over } u} Y
	\]
	where:
	\begin{itemize}
		\item \( e \to Y \) is a cartesian morphism over \( u : I \to J \)
		\item \( e \to X \) is a vertical morphism over \( \mathrm{id}_I \)
	\end{itemize}
	
	\section*{Diagram}
	
	\[
	\begin{tikzcd}
		& e \arrow[dl, "\text{vertical}"'] \arrow[dr, "\text{cartesian over } u"] & \\
		X \arrow[d, "P"'] & & Y \arrow[d, "P"] \\
		I & & J \arrow[ll, "u"]
	\end{tikzcd}
	\]
	
	\section*{Why This Encodes a Morphism in the Opposite Direction}
	
	In the original fibration \( P \), morphisms go forward along base change: pulling back \( Y \) along \( u : I \to J \) yields \( u^*Y \in \mathbf{E}_I \).
	
	In \( P^{\mathrm{op}} \), we want a morphism from \( Y \in \mathbf{E}_J \) to \( X \in \mathbf{E}_I \) over \( u : J \to I \). Since we cannot directly go backward, we simulate this by:
	\begin{enumerate}
		\item Pulling back \( Y \) along \( u \) to get \( e = u^*Y \)
		\item Connecting \( e \) to \( X \) via a vertical morphism
	\end{enumerate}
	
	This span represents a morphism in \( \mathbf{E}^{\mathrm{op}} \) from \( Y \) to \( X \) over \( u \).
	
	\section*{Equivalence Classes Under Vertical Isomorphism}
	
	Spans are considered equivalent if there exists a vertical isomorphism \( v : e \to e' \) such that the following diagram commutes:
	\[
	\begin{tikzcd}
		& e \arrow[dl, "\text{vertical}"'] \arrow[dr, "\text{cartesian}"] \arrow[dd, "v", "\cong"'] & \\
		X & & Y \\
		& e' \arrow[ul, "\text{vertical}"] \arrow[ur, "\text{cartesian}"'] &
	\end{tikzcd}
	\]
	
	This equivalence is justified by the universal property of cartesian morphisms: any two cartesian lifts of the same morphism are uniquely isomorphic via a vertical morphism.
	
	\section*{Why \( P^{\mathrm{op}} \) Is a Category}
	
	By taking equivalence classes of spans under vertical isomorphism:
	\begin{itemize}
		\item Composition becomes strictly associative
		\item Identity morphisms are well-defined
		\item We avoid coherence data required in bicategories
	\end{itemize}
	
	Thus, \( P^{\mathrm{op}} \) forms a genuine category, not just a bicategory of spans.
	
\end{document}